\section{Introduction}\label{sec:introduction}

Statistical offices and central banks increasingly measure prices from retailers' websites
and scanner data \citep{cavallo2016bpp, eurostat2017scanner, eurostat2020webscraping}. These
sources carry millions of products, but a product enters a consumer price index only after it
has been assigned to an item of the official basket. Done by hand, that step does not scale;
done by a model, it introduces an error. The literature on classification for price
statistics reports this error as accuracy \citep{martindale2020semi, harms2019cbs,
bis2026spectrum}. A central bank does not target accuracy. It needs to know how far measured
inflation would move.

This paper puts that question to one language and one index. We classify Thai product text
into the in-scope items of the Thai CPI with four families of method: multilingual text
embeddings with simple classifiers, a fine-tuned multilingual encoder, few-shot prompting of
two commercial LLMs, and cascades that send a cheap classifier's unsure products to an LLM.
We score them on \NRefSpecs{} held-out official price-collection specifications, published
by the Ministry of Commerce with their official item codes, and weight every score by CPI
expenditure. We then pass each method's confusion matrix through a simulator of the Thai CPI,
built from public weights and item indices, which returns the error the classifier would add
to measured headline, food and core inflation in basis points. Three questions organise the
paper.

\emph{RQ1: how far does transfer from other languages go, and what are official labels
worth?} Thai labelled product data are scarce, while Open Food Facts holds millions of
products in other languages. An \NameEfive{} classifier trained only on non-Thai products
classifies \AccWEfiveMLPZeroShot\% of the evaluated CPI weight correctly. Up to
\CfgBudgetMax{} official labels add little, and all \NMocTrainval{} of them raise it to
\AccWEfiveMLP\%.

\emph{RQ2: what does accuracy cost?} On official Thai text the LLMs are far ahead:
\NameOpus{} reaches \AccWOpus\% and \NameHaiku{} \AccWHaiku\%, against \AccWXlmrFT\% for the
best cheap method. A cascade from the \NameEfive{} classifier to \NameOpus{} matches
\NameOpus{} at \CostPctCasEfiveMLPOpus\% of its cost. Most of the remaining LLM error comes from
two places where our description of the items differed from the office's practice, not from
misreading Thai.

\emph{RQ3: what does classification error cost in basis points?} With \NameOpus's confusion
matrix, headline year-on-year inflation would be misstated by \BpYoyRmseHeadlineOpus{} basis
points (RMSE; \CfgCI\% interval \BpYoyRmseHeadlineOpusLo--\BpYoyRmseHeadlineOpusHi) and food
inflation by \BpYoyRmseFoodOpus. With the \NameEfive{} classifier the errors are
\BpYoyRmseHeadlineEfiveMLP{} and \BpYoyRmseFoodEfiveMLP{} basis points. Accuracy does not rank
methods by this error: two classifiers with the same weighted accuracy differ several-fold in
basis points, because what matters is whether the confused items' prices moved apart.

The paper makes three contributions. First, it turns a confusion matrix into basis points of
measured inflation, with a decomposition that attributes the error to individual confusions
and an uncertainty band that reflects how few observations some rows rest on. The simulator
needs only public weights and item indices, so any office can apply it to its own
classifier. Second, it evaluates classification for a CPI in a non-European language, on
the office's own held-out specifications and with expenditure weights on both accuracy and
coverage. Third, the whole pipeline is public and reproducible: it uses only open data, every
number in the paper is generated by code, and the LLM runs cost USD \LLMSpendUSD{} in total.

The closest work is the BIS Innovation Hub's Project Spectrum \citep{bis2026spectrum}, which
classifies euro-area web prices and names testing other languages and back-testing indices as
next steps \citep[p.~30]{bis2026spectrum}. Our taxonomy, language, product universe and
labels differ from Spectrum's, so our numbers are not directly comparable with theirs.
Section~\ref{sec:related} reviews the literature, Section~\ref{sec:data} describes the data
and Section~\ref{sec:methods} the methods. Section~\ref{sec:results} answers the three
questions, Section~\ref{sec:discussion} draws lessons for central banks in Asia,
Section~\ref{sec:limitations} lists the limitations and Section~\ref{sec:conclusion}
concludes.
